\chapter{Options reference}\label{ch:options}

This chapter lists the options of \nwsaom{} by topic. The definition of each effect (what is counted, its RSiena name and its validation) is in Chapter~\ref{ch:effects} and Appendix~\ref{app:effects}; here you find which option selects it, which arguments it takes and what it can be combined with. The text is derived from the help files \texttt{nwsaom.sthlp} and \texttt{nwsaom\_remarks.sthlp}; if the two disagree with the program, the program is right and the help file should be reported. Options marked \emph{native} exist only in the C plugin: with a plugin that is too old, or with the Mata engine, they are refused with an error message (Chapter~\ref{ch:technical}).

\section{Syntax}

\begin{lstlisting}
nwsaom, wave1(netname) wave2(netname)  |  waves(namelist)
        outdegree [reciprocity]  |  outdegreeendow outdegreecreation [...]
        [ endow(effects) creation(effects) ]
        [ covariate, structural, interaction and alias options ]
        [ behavior(...) and its effects ]
        [ composition change and missing data options ]
        [ ratecov(varlist) ratecovcoef(numlist) nocenter ]
        [ symmetric symtype(string) ]
        [ unconditional rate0(numlist) theta0(numlist) k0(#) k3(#) firstg(#) seed(#) cores(#) ]

nwsaom multinet, net1(namelist) net2(namelist) [ eff1() eff2() cross1() cross2() ... ]
estat gof | estat mems
\end{lstlisting}
Every model needs the waves and the baseline effect \opt{outdegree}. Everything else is optional. Options that take a \emph{varlist} are repeated per variable: one effect with its own coefficient per variable, named \emph{effect}\texttt{\_}\emph{variable} in the spelling you used (for example \opt{samex(smoke1 sport1)} gives \texttt{samex\_smoke1} and \texttt{samex\_sport1}). In \texttt{e(b)}, and thus in \opt{theta0()}, the effect types appear in a fixed order and the variables within a type in the order listed; \texttt{matrix list e(b)} shows it.

\section{Observed waves}

\begin{longtable}{@{}>{\raggedright\arraybackslash}p{4.2cm}>{\raggedright\arraybackslash}p{10.8cm}@{}}
\toprule Option & Meaning \\ \midrule \endhead
\opt{wave1(netname)}, \opt{wave2(netname)} & The two waves of a model with exactly two waves. Both are required together; not combinable with \opt{waves()}. \\
\opt{waves(namelist)} & Two or more waves in temporal order. The $W-1$ periods are chained into one pooled fit with one rate per period. \\
\opt{tvcov(groups)} & Time-varying covariates (RSiena's \texttt{varCovar}). \texttt{tvcov(a1 a2 a3 | s1 s2 s3)} declares covariates with one variable per wave; \texttt{|} separates covariates. Period $p$ uses the value of wave $p$ (the value of the last wave is not used). A covariate is named after its first variable without trailing digits (here \texttt{a} and \texttt{s}) and used under that name in \opt{nodeocov()}, \opt{nodeicov()}, \opt{nodecov()}, \opt{nodematch()}, \opt{simcov()} or their aliases, for example \texttt{tvcov(a1 a2 a3) egox(a) simx(a)}. Centred by the mean over all periods. No missing values. Possible with \opt{behavior()} (network covariate effects only) and in two-mode networks; not with \opt{interact()}, \opt{symmetric} or \texttt{estat gof}. \\
\opt{actormode(\#)} & Two-mode networks: the mode whose nodes are the actors, 1 or 2; the default 2 is the rows of the affiliation matrix given to \texttt{nwset, bipartite}. \\
\bottomrule
\end{longtable}

\section{Baseline network effects}

\begin{longtable}{@{}>{\raggedright\arraybackslash}p{4.2cm}>{\raggedright\arraybackslash}p{10.8cm}@{}}
\toprule Option & Meaning \\ \midrule \endhead
\opt{outdegree} & The density effect, evaluation-function role. Required in every model unless the endowment/creation form is used. \\
\opt{reciprocity} & Reciprocated ties. \\
\opt{endow(effects)}, \opt{creation(effects)} & Endowment and creation roles of effects (RSiena effect types \texttt{endow} and \texttt{creation}). The list may contain \texttt{outdegree reciprocity transties cycle4 gwesp balance transtrip cycle3 transmedtrip transrectrip inpop outpop inact outactsqrt indegpopularity outpopularity outactivity inactivity} for the network and \texttt{linear quadratic avalt avsim} for a behavior. An effect named in both lists is split into its two roles and replaces the plain effect; naming it in only one list, or also giving the plain effect, is an error. For \texttt{gwesp} the decay is taken from \opt{gwesp(\#)}. \\
\opt{endow(eff:var)}, \opt{creation(eff:var)} & Covariate effects are split per variable: \texttt{endow(altx:alc) creation(altx:alc)}, where the effect is \texttt{egox}, \texttt{altx}, \texttt{samex}, \texttt{simx} or \texttt{nodecov}. Coefficients \texttt{altxendow\_alc} and \texttt{altxcreation\_alc}. No missing values yet. \\
\opt{outdegreeendow}, \opt{outdegreecreation} & The long forms for the outdegree effect. They are currently \emph{refused}: the two statistics (lost and gained ties) add up to the distance statistic of the rate, so with the rate estimated the model is not identified (RSiena reports a singular covariance matrix for it as well). \\
\opt{reciprocityendow}, \opt{reciprocitycreation} & Reciprocity split, used together with \opt{outdegree}. \\
\opt{transtiesendow}, \opt{transtiescreation} & \opt{transties} split. Not yet with \opt{symmetric}. \\
\bottomrule
\end{longtable}
The pairs \opt{outdegreeendow}/\opt{outdegreecreation}, \opt{reciprocityendow}/\ldots, \opt{transtiesendow}/\ldots and \opt{linearendow}/\ldots are the same thing as the \opt{endow()}/\opt{creation()} lists and still work. Endowment and creation effects work with \opt{waves()}, \opt{present()}, \opt{missnet()} and \opt{behavior()}; with \opt{structural()} they are rejected.

\section{Covariate effects}

Covariates are read from the current dataset by row. The options below take a varlist.

\begin{longtable}{@{}>{\raggedright\arraybackslash}p{4.2cm}>{\raggedright\arraybackslash}p{10.8cm}@{}}
\toprule Option (alias) & Meaning \\ \midrule \endhead
\opt{nodeocov()} (\opt{egox()}) & Ego (sender) covariate effect, RSiena's \texttt{egoX}. \\
\opt{nodeicov()} (\opt{altx()}) & Alter (receiver) covariate effect, \texttt{altX}. \\
\opt{nodematch()} (\opt{samex()}) & Homophily on a categorical attribute, exact match, \texttt{sameX}. \\
\opt{simcov()} (\opt{simx()}) & Similarity, $1-|x_i-x_j|/\text{range}$ minus its mean over all pairs, \texttt{simX}. \\
\opt{nodecov()} & Sum of the sender's and the receiver's own value (continuous covariate). \\
\opt{absdiffx()}, \opt{egoxaltx()}, \opt{diffx()}, \opt{higher()}, \opt{egosqx()}, \opt{altsqx()}, \opt{samexrecip()}, \opt{simrecipx()} & RSiena's \texttt{absDiffX}, \texttt{egoXaltX}, \texttt{diffX}, \texttt{higher}, \texttt{egoSqX}, \texttt{altSqX}, \texttt{sameXRecip}, \texttt{simRecipX}. \\
\opt{egoplusaltsqx()}, \opt{diffsqx()}, \opt{egodiffx()}, \opt{althigheregox()}, \opt{unequalx()}, \opt{egolthresholdx()}, \opt{egorthresholdx()}, \opt{altlthresholdx()}, \opt{altrthresholdx()}, \opt{degabsdiffx()}, \opt{degposdiffx()}, \opt{degnegdiffx()}, \opt{covthreshold(\#)} & Covariate contrast and threshold effects. The four threshold effects compare the centred covariate with \opt{covthreshold(\#)} (default 0). Directed relations, no missing values. Native. \\
\opt{dyadcov(matnames)}, \opt{dyadcovrecip(matnames)} & Dyadic covariate (RSiena's \texttt{X} and \texttt{XRecip}). Each name is a Stata $n\times n$ matrix in the node order of the waves; it is centred by the mean of its non-missing off-diagonal cells (not with \opt{nocenter}), missing cells count as 0 after centring. One coefficient \texttt{dyadcov\_}\emph{name} per matrix. Not with \opt{interact()}. \\
\opt{tvdyadcov(stem \ldots)}, \opt{tvdyadcovrecip(stem \ldots)} & Changing dyadic covariates: matrices \emph{stem}1 \ldots{} \emph{stem}$(W-1)$, one per period. \\
\opt{nocenter} & Use the raw values of the covariates of \opt{egox()}, \opt{altx()}, \opt{nodecov()} and \opt{ratecov()}. By default they are centred by their means, as RSiena's \texttt{coCovar(centered = TRUE)}; missing values are set to the mean. The means are returned in \texttt{e(covmeans)}. \opt{samex()} does not depend on centring; \opt{simx()} is centred by its similarity mean either way. \\
\bottomrule
\end{longtable}
Missing covariate values: for \opt{absdiffx()}, \opt{egoxaltx()}, \opt{altx()}, \opt{egox()}, \opt{samex()} and \opt{simx()} a missing value is imputed with the mean when actors decide, but the statistic skips the ties whose relevant actor has a missing value, as in RSiena.

\section{Structural network effects}

\begin{longtable}{@{}>{\raggedright\arraybackslash}p{4.2cm}>{\raggedright\arraybackslash}p{10.8cm}@{}}
\toprule Option & Meaning \\ \midrule \endhead
\opt{indegpopularity}, \opt{outpopularity}, \opt{outactivity}, \opt{inactivity} & Square-root and squared degree effects (\texttt{inPopSqrt}, \texttt{outPopSqrt}, \texttt{outAct}, \texttt{inActSqrt}). \\
\opt{inpop}, \opt{outpop}, \opt{inact}, \opt{outactsqrt} & RSiena's \texttt{inPop}, \texttt{outPop}, \texttt{inAct}, \texttt{outActSqrt}. \opt{outpop} and \opt{inact} have the same total statistic and cannot be used together. \\
\opt{inpopc}, \opt{outpopc}, \opt{outactc}, \opt{inactc}, \opt{degplus}, \opt{degplusc}, \opt{avdeg(\#)}, \opt{transtriads} & Centred and combined degree effects, the average-degree effect and the triangle count; \opt{degplus}, \opt{degplusc} and \opt{transtriads} for non-directed relations only. Native. \\
\opt{outtrunc(\#)}, \opt{outmore(\#)}, \opt{outthreshold(\#)}, \opt{outinv(\#)}, \opt{outsqinv(\#)} & Functions of the outdegree with a parameter; one term per value given. \\
\opt{transtrip}, \opt{transties}, \opt{transmedtrip}, \opt{transrectrip}, \opt{cycle3}, \opt{cycle4}, \opt{balance}, \opt{sqrtcycle4}, \opt{between} & Triadic and cyclic effects. \opt{balance} has no parameter (its constant is computed from the data). \\
\opt{gwesp(\#)} & Geometrically weighted edgewise shared partners, the direct decay value (RSiena's parameter divided by 100). \\
\opt{gwespfb(\#)}, \opt{gwespbf(\#)}, \opt{gwespbb(\#)}, \opt{gwesprr(\#)}, \opt{gwdspff(\#)}, \opt{gwdspfb(\#)} & The other weighted shared-partner versions, with the same decay scale. \\
\opt{nbrdist2}, \opt{dblnbrdist2} & Actors at distance two (once, at least twice). \\
\opt{recipop}, \opt{recipopsqrt}, \opt{reciact}, \opt{reciactsqrt} & Reciprocal-degree effects. \\
\opt{isolatenet}, \opt{outiso}, \opt{antiiso}, \opt{antiiniso}, \opt{antiiniso2}, \opt{inplus3}, \opt{isolatepop} & Isolate effects. They are weakly identified on small or sparse networks; \opt{antiiso} and \opt{isolatepop} can destabilise the estimation there. \\
\opt{outoutass}, \opt{ininass}, \opt{outinass}, \opt{inoutass} & Degree assortativity, default parameterisation. \\
\opt{ownx(effect[:\#] \ldots)} & Further own-network effects without a covariate (native), among them \texttt{jin jout transtrip1 transtrip2 transrectrip2 transtripfr transtripfe transtripee}, the \texttt{outpopmore}, \texttt{outpopsqrtmore}, \texttt{outpopthreshold} families, the ego degree effects, the assortativity effects with parameter 2 and the primary-setting effects. \\
\opt{netx(effect:var \ldots)}, \opt{xparam(\#)} & Further covariate network effects (native), for example \texttt{crossxoutact}, \texttt{homxoutact2}, \texttt{transtripx}, \texttt{altdist2}, \texttt{sameinxcycle4}, \texttt{samexv:v1:v2}; \opt{xparam()} is the parameter of some of them. \\
\opt{samextranstrip()}, \opt{diffxtranstrip()}, \opt{simxtranstrip()}, \opt{homxtranstrip()}, \opt{samexoutact()}, \opt{diffxoutact()}, \opt{samexreciact()}, \opt{diffxreciact()}, \opt{homxoutact()}, \opt{altxoutact()}, \opt{inpopx()}, \opt{outpopx()}, \opt{inactx()}, \opt{outactx()}, \opt{samexinpop()}, \opt{diffxinpop()} & Covariate-by-triad, popularity and activity effects. One numeric covariate per name, directed one-mode networks, no missing values. Native. \\
\opt{dyadx(effect:matrix \ldots)} & Network effects of a constant dyadic covariate (\texttt{wwx}, \texttt{cywwx}, \texttt{inwwx}, \texttt{outwwx}, \texttt{wxx}, \texttt{xwx}, \texttt{xwx1}, \texttt{xwx2}, \texttt{xxw}). Matrix used raw. Native. \\
\opt{tvdyadx(effect:stem \ldots)} & The same with changing dyadic covariates. \\
\bottomrule
\end{longtable}
{\sloppy The structural effects that are meaningless in a non-directed network (\opt{reciprocity}, \opt{cycle3}, \opt{inactivity}, \opt{outpopularity}, \opt{ininass}, \opt{inoutass}, \opt{outoutass}, \opt{antiiso}, \opt{isolatepop}, \opt{transrectrip}, \opt{transtrip}) are rejected when the relation is non-directed. In a two-mode network the available effects are \opt{outdegree}, \opt{indegpopularity}, \opt{outactivity}, \opt{cycle4}, \opt{outiso}, \opt{antiiniso}, \opt{antiiniso2}, \opt{inplus3}, \opt{gwdspfb()}, \opt{avdeg()}, \opt{outtrunc()}/\opt{outmore()}/\opt{outthreshold()}, \opt{ownx()}, and the covariate effects \opt{nodeocov()} (on the actors), \opt{nodeicov()} (on the events), \opt{egosqx()}, \opt{altsqx()} and some \opt{netx()} effects.\par}

\section{Interactions}

\begin{longtable}{@{}>{\raggedright\arraybackslash}p{4.2cm}>{\raggedright\arraybackslash}p{10.8cm}@{}}
\toprule Option & Meaning \\ \midrule \endhead
\opt{interact(e1\#e2[\#e3] \ldots)} & Two- or three-way interaction of effects already in the model (RSiena's \texttt{includeInteraction}), with its own coefficient. Several interactions may be listed. RSiena's rule applies: effects are ego effects, dyadic effects or other effects, and two effects need an ego effect or two dyadic effects, three effects two ego effects or only ego and dyadic effects. A covariate effect for several variables is named by its coefficient name, for example \texttt{interact(samex\_smoke1\#egox\_alcohol1)}; with one variable the effect type alone works. Coefficients are named \texttt{interact\_}\emph{A}\texttt{\_}\emph{B} (\texttt{ix\_}\emph{A}\texttt{\_}\emph{B} if longer than 32 characters). Behavior interactions are not supported. \\
\bottomrule
\end{longtable}

\section{Behavior co-evolution}

\begin{longtable}{@{}>{\raggedright\arraybackslash}p{4.2cm}>{\raggedright\arraybackslash}p{10.8cm}@{}}
\toprule Option & Meaning \\ \midrule \endhead
\opt{behavior(varlist)} & A bounded integer behavior, one variable per wave in temporal order. Several behaviors: \texttt{behavior(a1 a2 a3 : linear avalt | s1 s2 s3 : linear)}, separated by \texttt{|}, each optionally followed by \texttt{:} and its own effects. Coefficients are named \texttt{beh\_}\emph{name}\texttt{\_}\emph{effect}. \\
\opt{linear} & Linear shape effect, the behavior's baseline; required with \opt{behavior()} unless \opt{linearendow}/\opt{linearcreation} are given. (The latter are currently refused, for the same reason as \opt{outdegreeendow}.) \\
\opt{quadratic}, \opt{avalt}, \opt{avsim} & Quadratic shape; average-alter influence; average-similarity influence. Each also in an endowment/creation form (\opt{quadraticendow} \opt{quadraticcreation}, \opt{avaltendow} \opt{avaltcreation}, \opt{avsimendow} \opt{avsimcreation}). \\
\opt{behsim}, \opt{behegox}, \opt{behaltx}, \opt{behegoplusalt} & Selection effects on the network side that use the current behavior (similarity, ego, alter, ego plus alter). Network coefficients; require \opt{behavior()}. \\
\opt{beheff(tokens)} & Behavior effects of network position, alters and similarity (the tokens are listed in the help file); some are native only. \\
\opt{behfrom(varlist)} & Covariate main effects on the behavior (\texttt{effFrom}); at most 8 covariates in total with \opt{behx()}. \\
\opt{behx(family:var \ldots)} & Behavior effects of the network times a covariate; \emph{var} may also be \texttt{beh2}, \texttt{beh3}, \ldots, another behavior. Native. \\
\opt{behavxalt()}, \opt{behavsamexalt()}, \opt{behtotxalt()}, \opt{behtotsamexalt()} & Covariate-moderated influence (\texttt{avXAlt}, \texttt{avSameXAlt}, \texttt{totXAlt}, \texttt{totSameXAlt}). \\
\opt{behw(family:matrix \ldots)}, \opt{tvbehw(family:stem \ldots)}, \opt{behwparam(\#)} & Behavior effects of the network times a constant or changing dyadic covariate (native). \\
\opt{behshape(effect[:\#] \ldots)}, \opt{behnet2(effect[:\#] \ldots)} & Further shape effects and further behavior effects of the network (native). \\
\opt{netbeh(effects)} & Network effects with the dependent behavior as covariate (native). First behavior only; coefficients \texttt{bv\_}\emph{effect}. \\
\opt{behtheta0(numlist)} & Starting values of the behavior-side coefficients. \\
\bottomrule
\end{longtable}

\section{Composition change, missing data and structural values}

\begin{longtable}{@{}>{\raggedright\arraybackslash}p{4.2cm}>{\raggedright\arraybackslash}p{10.8cm}@{}}
\toprule Option & Meaning \\ \midrule \endhead
\opt{present(varlist)} & Actors present at each wave: one 0/1 variable per wave. An actor is present in a period if it is present at both of its waves. \\
\opt{joinat(varname)}, \opt{leaveat(varname)} & Joiners and leavers at fractional times within a period, exactly as RSiena's \texttt{sienaCompositionChange()} option 1; time scale wave 1 = 1, last wave = $W$. A missing value means from the first wave, respectively to the last. One interval per actor. \\
\opt{missnet(matlist)} & One 0/1 $n\times n$ matrix per wave marking missing dyads. \\
\opt{missbeh(varlist)} & One 0/1 variable per wave marking missing behavior values; requires \opt{behavior()}. \\
\opt{structural(matname)} & One 0/1 matrix marking dyads whose value is fixed by design. Exactly two waves (\opt{wave1()}/\opt{wave2()}, not \opt{waves()}), network only; not with \opt{symmetric}, \opt{ratecov()} or the network endowment/creation split. \\
\bottomrule
\end{longtable}
Composition change implies unconditional estimation. Data preparation (coding of absent actors' ties and of the missing values) is the user's responsibility; see the chapter on composition change and missing data.

\section{Rates}

\begin{longtable}{@{}>{\raggedright\arraybackslash}p{4.2cm}>{\raggedright\arraybackslash}p{10.8cm}@{}}
\toprule Option & Meaning \\ \midrule \endhead
\opt{ratecov(varlist)} & Covariates that change the rate: actor $i$'s rate is the period rate times $\exp(\sum b_k x_k[i])$. One coefficient per variable, shared by all periods. Works with \opt{waves()} and \opt{symmetric}. Not with co-evolution. \\
\opt{ratecovcoef(numlist)} & Starting values of these coefficients, one per variable or one for all. Default 0. \\
\opt{outrate}, \opt{inrate}, \opt{reciprate} and their \texttt{inv}/\texttt{log} versions (\opt{outrateinv}, \opt{outratelog}, \ldots) & Structural rate effects: the rate depends on the actor's current out-, in- or reciprocated degree (or its inverse or logarithm). Each adds a covariate-rate coefficient. Directed one-mode relations, plain waves only. Native. \\
\bottomrule
\end{longtable}

\section{Non-directed relations}

\begin{longtable}{@{}>{\raggedright\arraybackslash}p{4.2cm}>{\raggedright\arraybackslash}p{10.8cm}@{}}
\toprule Option & Meaning \\ \midrule \endhead
\opt{symmetric} & Model a non-directed relation. It is not needed when the waves are declared undirected (\texttt{nwset, undirected}) or when all directed waves happen to be tie-symmetric: \nwsaom{} then models a non-directed relation automatically, as RSiena does. A tie-symmetric directed wave stays directed with \opt{reciprocity} or with effects for which the non-directed model does not apply. With directed waves that are not tie-symmetric \opt{symmetric} is an error: it does not symmetrise data. Combinable with \opt{waves()}, \opt{present()}, \opt{missnet()}, \opt{ratecov()} (network only) and co-evolution (native plugin required). \\
\opt{symtype(string)} & The model type, by name or by RSiena's \texttt{modelType} number: \texttt{forcing} (2, the default), \texttt{confirmation} (3), \texttt{force} (4), \texttt{agree} (5), \texttt{joint} (6). \\
\bottomrule
\end{longtable}

\section{Estimation control}

\begin{longtable}{@{}>{\raggedright\arraybackslash}p{4.2cm}>{\raggedright\arraybackslash}p{10.8cm}@{}}
\toprule Option & Meaning \\ \midrule \endhead
\opt{unconditional} & Unconditional method of moments for a network-only model; the rates become method-of-moments parameters. The default for a network-only model is conditional estimation. Co-evolution, \texttt{nwsaom multinet} and \texttt{nwsaom multiplex} models are always unconditional. \\
\opt{rate0(numlist)} & Starting value(s) of the rate(s) of an unconditional network-only fit: one value or one per period. Default: RSiena's closed-form value from the data. Not used under conditional estimation. \\
\opt{theta0(numlist)} & Starting values of the coefficients, one per effect in the order of \texttt{e(b)} (the error message lists the order if the argument is missing or has the wrong length). Default all zero, except that \opt{outdegree} of a network-only model starts at RSiena's data-derived value. \\
\opt{k0(\#)} & Number of phase-1 replicates, default 50. \\
\opt{k3(\#)} & Number of phase-3 replicates, default 1\,000. \\
\opt{firstg(\#)} & Initial gain of phase 2, default 0.2. \\
\opt{seed(\#)} & Random-number seed. \\
\opt{cores(\#)} & Threads of the native simulator; 0 (default) is all physical cores, 1 is single-threaded. Results for a given seed do not depend on it. \\
\bottomrule
\end{longtable}
{\sloppy Fragments of this chapter that were run with Stata (three waves of the Glasgow data, \texttt{k0(20) k3(50) seed(7)} added to shorten the runs) and finished without error: \texttt{outdegree reciprocity endow(transties) creation(transties)}; \texttt{outdegree reciprocity ratecov(alcohol1)}; \texttt{outdegree reciprocity tvcov(alcohol1 alcohol2 alcohol3) egox(alcohol)}; \texttt{samex(smoke1) egox(alcohol1) interact(samex\_smoke1\#egox\_alcohol1)}; \texttt{outdegree reciprocity egox(alcohol1) nocenter}; and \texttt{outdegree reciprocity gwesp(0.69) unconditional cores(1)}.\par}

\section{Several dependent networks: \texttt{nwsaom multinet}}

\begin{lstlisting}
nwsaom multinet, net1(namelist) net2(namelist) [eff1(effects) eff2(effects)
        cross1(effects) cross2(effects) sym1 sym2 k0(#) k3(#) firstg(#) seed(#)]
\end{lstlisting}
Two directed one-mode networks on the same actors evolve together (RSiena's two dependent network variables); estimation is unconditional, with a rate per network and period. Native only.
\begin{itemize}[nosep]
  \item \opt{net1()} and \opt{net2()} list the waves of the two networks; the number of waves is the same, at least two.
  \item \opt{eff1()} and \opt{eff2()} give the within-network effects of each network; default \texttt{outdegree reciprocity}. Possible are \texttt{outdegree reciprocity transtrip cycle3 transmedtrip cycle4 indegpop outact} and, with a variable, \texttt{egox:}\emph{var}, \texttt{altx:}\emph{var}, \texttt{samex:}\emph{var}, \texttt{simx:}\emph{var}.
  \item \opt{cross1()} lists the cross-network effects of network 1 that depend on network 2, \opt{cross2()} those of network 2 on network 1 (\texttt{crprod crprodrecip crprodmutual from frommutual fromany to toback torecip clxwx clxwx1 clxwx2 mixedinxw}, and many further families with and without a covariate; a parameter follows a colon, for example \texttt{from:2}). The cross-network statistics read the other network at the start of the period, the simulation its current state, as RSiena does.
  \item \opt{sym1} and \opt{sym2} declare the first, respectively second network non-directed.
\end{itemize}

\section{Postestimation commands}

\begin{description}[nosep,style=nextline]
  \item[\texttt{estat gof}] Goodness of fit with RSiena's Mahalanobis-distance test for the outdegree, indegree and geodesic distributions. Options and the returned \texttt{r()} results are in \texttt{help nwsaom\_estat}.
  \item[\texttt{estat mems}] Micro effects on macro structure sensitivity analysis (Duxbury); \texttt{help nwsaom\_estat}.
  \item[\texttt{nwsaom\_meta}] Combines the same model fitted to several networks, the counterpart of RSiena's \texttt{siena08}; \texttt{help nwsaom\_meta}.
\end{description}

\section{Combinations that are not available}\label{sec:combinations}

The following list collects the restrictions that the program enforces and that the help file states; a combination that is not allowed stops with an error message and never gives a silently different model.
\begin{itemize}[nosep]
  \item {\sloppy \opt{wave1()}/\opt{wave2()} together with \opt{waves()}; \opt{outdegree} together with \opt{outdegreeendow}/\opt{outdegreecreation}, or any plain effect together with its own endowment/creation split.\par}
  \item \opt{structural()}: two waves only, no \opt{behavior()}, no \opt{symmetric}, \opt{ratecov()} or network endowment/creation.
  \item \opt{ratecov()}: not with co-evolution; with \opt{symmetric} not together with \opt{present()}, \opt{joinat()} or \opt{leaveat()}, and not with co-evolution.
  \item \opt{tvcov()}: no missing values; not with \opt{interact()}, \opt{symmetric} or \texttt{estat gof}.
  \item Several behaviors: not with \opt{symmetric}, \opt{ratecov()}, \opt{present()}, \opt{missnet()}/\opt{missbeh()}, \opt{structural()}, network endowment/creation, or \texttt{estat gof}/\texttt{estat mems}.
  \item \opt{joinat()}/\opt{leaveat()}: not with \opt{behsim} or endowment/creation effects in co-evolution models; not yet with \opt{endow()}/\opt{creation()}; with \opt{symmetric} not with \opt{ratecov()}.
  \item Valued networks and networks carrying \texttt{nwset}-style temporal metadata are rejected.
\end{itemize}
